\begin{frame}
\frametitle{Ukázka: QtRvSim vect-inc, vect-add2, vect-add}

Začneme s přímomapovanou skrytou pamětí (direct mapped cache) a \texttt{vect-inc}. Při velikosti bloku větší než jedno slovo a velikosti vektoru větší než kapacita jedné cesty bude u \texttt{vect-add2} docházet k značné degradaci výkonu. Po navýšení stupně asociativity na dvě cesty i při zachování celkové kapacity již ke kolizím při každé iteraci docházet nebude. Ovšem po rozšíření algoritmu na práci s třemi vektory (\texttt{vect-add}) opět problém při zpětném/odloženém zápisu (write-back) nastane a nejhorší situace nastane pro politiku záměny nejdéle nepoužívaného prvku (LRU). Zvýšení stupně na tři, nebo spíše obvyklejší čtyři, cesty se opět pro na začátek prázdnou cache počet výpadků omezí na jeden za každý načatý blok.

\begin{itemize}
\item \url{https://gitlab.fel.cvut.cz/b35apo/stud-support/-/tree/master/seminaries/qtrvsim/vect-inc}
\item \url{https://gitlab.fel.cvut.cz/b35apo/stud-support/-/tree/master/seminaries/qtrvsim/vect-add2}
\item \url{https://gitlab.fel.cvut.cz/b35apo/stud-support/-/tree/master/seminaries/qtrvsim/vect-add}
\end{itemize}

\end{frame}
